\documentclass[11pt]{article}
\usepackage[margin=1in]{geometry}
\usepackage[T1]{fontenc}
\usepackage[utf8]{inputenc}
\usepackage{amsmath,amssymb,amsthm}
\usepackage{enumitem}

\title{ICPC World Finals 2024\\C. Citizenship}
\author{}
\date{}

\begin{document}
\maketitle

\section*{Problem Summary}

For an application date $A$, the previous $y$ years are the $y$ consecutive intervals
\[
[A-365, A-1], [A-730, A-366], \dots, [A-365y, A-365(y-1)-1].
\]
In each such interval we must have been present at least $d$ days, which is equivalent to being absent at
most $365-d$ days. We must output the earliest valid application date after the last recorded trip.

\section*{Key Observations}

\begin{itemize}[leftmargin=*]
    \item The only thing that matters about a date is its position modulo $365$. For a fixed residue
    $r = A \bmod 365$, the relevant 12-month periods are aligned to a fixed partition of the timeline into
    blocks of length $365$.
    \item Therefore there are only $365$ possible alignments to check.
    \item For one fixed residue $r$, define block $b$ as
    \[
        [r + 365b,\; r + 365b + 364].
    \]
    The application date $A = r + 365t$ is valid exactly when blocks
    $t-y, t-y+1, \dots, t-1$ are all good, where a block is good if it contains at most $365-d$ absent days.
    \item The absence intervals are few, but the total date range is still small enough: at most about
    $5000 \cdot 365 + 1001 \cdot 365$. A plain prefix-sum over days is easily fast enough.
\end{itemize}

\section*{Algorithm}

\begin{enumerate}[leftmargin=*]
    \item Convert every calendar date to an integer day number using the fixed 365-day calendar.
    \item Mark all absent days with a difference array and build a prefix sum of absent days.
    \item Let $L$ be the last day in the input. It is enough to search up to
    $L + 365(y+1)$, because after that all $y$ required years lie completely after the last trip and are
    therefore automatically valid.
    \item For each residue $r \in \{0,\dots,364\}$:
    \begin{enumerate}[leftmargin=*]
        \item Scan the 365-day blocks with that alignment.
        \item For each block, compute its absent-day count in $O(1)$ from the prefix sum and decide
        whether the block is good.
        \item Maintain the current run length of consecutive good blocks.
        \item As soon as the run length reaches $y$, the application date immediately after the last block in
        that run is valid. If it is after $L$, update the global minimum answer.
    \end{enumerate}
    \item Convert the minimum day number back to \texttt{YYYY-MM-DD}.
\end{enumerate}

\section*{Correctness Proof}

We prove that the algorithm outputs the earliest valid application date.

\paragraph{Lemma 1.}
For a fixed residue $r$, an application date $A \equiv r \pmod{365}$ is valid if and only if the $y$
consecutive 365-day blocks immediately before $A$ in the $r$-aligned partition are all good.

\paragraph{Proof.}
Write $A = r + 365t$. Then
\[
[A-365, A-1] = [r + 365(t-1),\; r + 365t - 1],
\]
which is exactly block $t-1$. Similarly, the next earlier year is block $t-2$, and so on down to block
$t-y$. So the $y$ periods required by the statement are exactly the $y$ consecutive blocks immediately
before $A$. Each period is valid exactly when its absent-day count is at most $365-d$. \qed

\paragraph{Lemma 2.}
For a fixed residue $r$, when the scan first finds a run of at least $y$ consecutive good blocks ending at
block $b$, the date $A = r + 365(b+1)$ is the earliest valid application date with residue $r$.

\paragraph{Proof.}
By Lemma 1, $A$ is valid because the preceding $y$ blocks are good. Any earlier date with the same
residue has the form $r + 365t$ with $t \le b$. Its preceding $y$ blocks must therefore end at some block
strictly before $b$. Since $b$ is the first point where the scan sees $y$ consecutive good blocks, no such
earlier date can be valid. \qed

\paragraph{Lemma 3.}
There exists a valid application date no later than $L + 365(y+1)$, where $L$ is the last day in the input.

\paragraph{Proof.}
Consider any date $A > L + 365y + 364$. Then the earliest of the $y$ required periods starts at
$A - 365y > L$. So all $y$ required 365-day periods lie completely after the last recorded trip and contain
zero absent days. Hence they are all valid. \qed

\paragraph{Theorem.}
The algorithm outputs the earliest valid application date after the last input date.

\paragraph{Proof.}
By Lemma 3, the search range always contains at least one valid date. For each residue class, Lemma 2
gives the earliest valid date with that residue. Taking the minimum over all 365 residues therefore yields
the earliest valid date overall. \qed

\section*{Complexity Analysis}

Let $T$ be the number of days scanned. Here
\[
T = O(5000 \cdot 365 + y \cdot 365),
\]
which is at most about $2.2 \cdot 10^6$.

Building the difference array and prefix sums takes $O(T + n)$. The 365 residue scans together also
process only $O(T)$ blocks. Thus the total running time is $O(T + n)$ and the memory usage is $O(T)$.

\section*{Implementation Notes}

\begin{itemize}[leftmargin=*]
    \item Since the calendar is fixed and has no leap years, date conversion is just base-365 arithmetic with
    the given month lengths.
    \item The output must be \emph{after} the last date in the input, so the candidate date must satisfy
    $A > L$.
    \item Difference arrays handle inclusive absence intervals cleanly: add $+1$ at the start day and $-1$
    at the day after the end.
\end{itemize}

\end{document}
